\documentclass[12pt]{article}
\usepackage[utf8]{inputenc}
\usepackage[T1]{fontenc}
\usepackage{hyperref}
\title{Bermejo Wich\'i Language Resources for K-12 Curriculum Development}
\author{Compiled by Claude Sonnet 5.5 (claude-sonnet-5-5) for Andr\'as Kornai}
\date{}
\begin{document}
\maketitle

\section{Introduction}
Bermejo Wich\'i (ISO 639-3 wlv; Glottocode wich1263) belongs to the Mataguayan (Matacoan) family and is spoken in northern Argentina. The searches report that Wich\'i as a whole has about 40,000 speakers in Argentina and Bolivia. They put the Bermejo variety at roughly 3,000 speakers in Formosa, though the survey's figure is about 25,000 for the code. Wich\'i is described as the second most spoken language in Argentina after Spanish. More material exists for Wich\'i generally than for the other languages in this batch.

\section{Primary Educational Resources}
\begin{itemize}
\item Since 2011, a CONICET-linked project has run workshops in the Bajo Bermejo region of Chaco. Teachers and community members produce books, posters and educational games for bilingual intercultural teaching. A related project, coordinated by Virginia Unamuno with Camilo Ballena of the CIFMA, supports teachers in producing materials. Illustrated bilingual Wich\'i-Spanish mathematics books were reported in 2026. The CONICET tag page lists news items on these: \url{https://www.conicet.gov.ar/tag/wichi/}
\end{itemize}
I did not find the individual textbooks online.

\section{Linguistic and Reference Materials}
\begin{itemize}
\item DIWICA, the online Wich\'i-Castellano dictionary (2021), led by Ver\'onica Nercesian at CONICET. It covers the Pilcomayo and Bermejo dialect complexes. A review in \emph{Revista de Lexicograf\'ia} gives its address as diccionariowichi.com.ar. That address did not appear as a search result, so I give it as text only. The review: \url{https://revistas.udc.gal/index.php/rlex/article/view/rlex.2022.28.1.9489/g9489_pdf}
\item Documents hosted on the Lengua Wich\'i site, with the file names ``tsalanawufichas'' (2006) and ``Vidal and Nercesian 2009''. I did not check their contents: \url{https://lenguawichi.com.ar/wp-content/uploads/2017/01/2006-tsalanawufichas.original.pdf} and \url{https://lenguawichi.com.ar/wp-content/uploads/2017/01/vidal-nercesian-2009.original.pdf}
\end{itemize}

\section{Religious and Community Resources}
None were found in the searches.

\section{Digital and Audio Resources}
DIWICA is the main digital resource. No audio collections were found in the searches.

\section{Curriculum Development Context}
Wich\'i has a long record of intercultural bilingual education work in Chaco, Formosa and Salta. This includes teacher-made materials and a 2021 Universidad Nacional de Salta diploma in Wich\'i-Castellano interpretation and translation. A curriculum can draw on the CONICET workshops and on DIWICA. Dialect differences between Bermejo and Pilcomayo varieties need care, and the work should be done with Wich\'i teachers.

\end{document}